\chapter{Összegzés}
\label{ch:sum}

Realisztikusnak minősülő vegetációk generálása nem egy egyszerű feladat, tekintve a növényzet által lefedett területek nagyságát, komplexitását, nem is beszélve a természetben található növényzeteket alkotó fajok diverzitásáról. Ezért is fontos a különböző elérhető adatok, illetve eszközök vizsgálata. Segítségükkel ugyanis a növényzeteket formáló hatásokat lehet könnyebben összefoglalni és szabályokat megfogalmazni, amelyeket aztán pontosabb eszközökben tudunk alkalmazni.  A szakdolgozatban bemutatott tájékozódási futó térképek is láthatóan egy ilyen jól kezelhető információ forrásnak bizonyulnak, amelyekkel végül nagyobb területeken elterülő vegetációkat lehet hatékonyan leírni. A különböző szabványok láthatólag ugyanis egy kellő keretet tudnak adni az ezeken térképeken használt objektumok leírására, amelyekkel így hatékony lehet különböző generálási szabályokat alkalmazni az adott régiókra. 
Továbbá az is meg lett mutatva, hogy az ezen adatforrások jól beillenek az már bevált növényzet szimulálási módszerekbe és technikákba, hiszen

Ennek az új adatforrásnak az értelmezésével, illetve korábban már jó bevált módszerek használatával pedig egy olyan szerelőszalag-szerű folyamatot lehetséges összeállítani, amellyel    


\todo{Összegzés}